\documentclass[11pt]{article}
\usepackage[margin=2.6cm]{geometry}
\usepackage{amsmath,amssymb,amsthm}
\usepackage{enumitem}
\usepackage{xspace}
\usepackage[T1]{fontenc}
\setlist{itemsep=1pt,topsep=3pt}

\newcommand{\sys}{\textsf{Unit}\xspace}
\newcommand{\kw}[1]{\texttt{#1}}
\newcommand{\U}{\mathsf{U}}
\newcommand{\Id}{\mathsf{Id}}
\newcommand{\PP}{\mathbb{P}^1}
\newcommand{\tr}[1]{\lVert #1\rVert}
\newcommand{\Spec}{\mathrm{Spec}}

\theoremstyle{remark}
\newtheorem{remark}{Remark}

\title{Synthetic algebraic geometry in \sys:\\what has been formalised}
\author{}
\date{21 October 2026}

% stamp on every page: file name and date/time of compilation
\newcount\stamphour \newcount\stampmin
\stamphour=\time \divide\stamphour by 60
\stampmin=\time \multiply\stamphour by 60 \advance\stampmin by -\stamphour \divide\stamphour by 60
\edef\stamptime{\the\year-\ifnum\month<10 0\fi\the\month-\ifnum\day<10 0\fi\the\day\space
 \ifnum\stamphour<10 0\fi\the\stamphour:\ifnum\stampmin<10 0\fi\the\stampmin}
\AddToHook{shipout/foreground}{\put(0.5\paperwidth,-\dimexpr\paperheight-0.45in\relax)%
 {\makebox[0pt]{\scriptsize\ttfamily\jobname.pdf\quad\stamptime}}}
\begin{document}
\maketitle

\begin{abstract}
\noindent
We report on a first formalisation of parts of \emph{A foundation for
synthetic algebraic geometry} (Cherubini, Coquand, Hutzler~\cite{SAG})
in the prototype \sys, the version where recursion is only through
eliminators. The foundation is a univalent universe with Rijke's axiom
of replacement, as formulated in~\cite{CavCoq}; propositional truncation
and set quotients are derived from it. Finitely presented algebras are
given by presentations and a universal property, and every presentation
has a presented algebra. SQC is one uniform hypothesis on all finitely
presented algebras. The results checked are Lemma~1.3.8,
Propositions~2.2.2--2.2.3, Lemma~2.2.4 (from Loc and SQC only),
Sections~3 and~4 (including
two proofs of Theorem~3.3.3), Section~5 under the full SAG
axioms (gluing has its universal property for small set targets),
Example~6.1.9 (the
fibres computed in the example, and the coordinate algebra of a fibre),
Lemma~6.2.1, Section~7.1 except Remark~7.1.3 (the localised module, 7.1.2 and 7.1.8 as
equivalences of modules, and 7.1.9--7.1.13 under the full SAG axioms),
Lemma~7.3.5, Theorem~7.3.6 (from the SAG axioms and fibrewise weak
quasi-coherence), Lemma~7.4.2, Theorem~7.4.7 and
Theorem~7.4.8. From \emph{Projective space in synthetic algebraic
geometry}~\cite{Proj}, Proposition~4.6 is checked: line bundles on
$\PP$ are $\mathsf{Line}\times\mathbb Z$. It uses the rank-one case of
Horrocks' theorem, proved for any commutative ring. Also checked are the
two definitions of $\mathbb P^n$ and their equivalence, Lemma~1.6,
Proposition~1.7, and Section~7 of~\cite{Proj}: line bundles on
$\mathbb P^n$, $n\ge 2$, are $\mathsf{Line}\times\mathbb Z$; and
Corollary~5.6: $\mathrm{Aut}(\mathbb P^n)=\mathrm{PGL}_{n+1}(R)$ for
$n\ge2$. The combined development through Section~5 checks in
$1.30$ seconds.
\end{abstract}

\section{The foundation}

All developments share one segment whose parameters are the
foundational axioms:
\begin{itemize}
\item \emph{univalence}, at every level, in the form ``$\mathsf{idtoeqv}$
  is an equivalence'';
\item \emph{replacement}, as two parameters which together form one
  axiom schema~\cite[Definitions 2, 4, 6]{CavCoq}:
  \begin{itemize}
  \item \kw{images}: for $A:\U_l$, a locally $\U_l$-small $B:\U_k$ and
    $f : A\to B$, a chosen factorisation
    $A \to I \to B$ with $I:\U_l$, a commuting triangle, and
    $I\to B$ an embedding;
  \item \kw{surjections}: the map $A\to I$ is surjective in the sense
    that precomposition is an equivalence on sections of every family of
    propositions over $I$, at every level $j$.
  \end{itemize}
\end{itemize}
The chosen image does not depend on the level $j$ of elimination. The
schema is split into data and a property only because a type quantifying
over all levels does not belong to a universe. No resizing is assumed.

From these parameters the following are derived, not assumed:
\begin{itemize}
\item function extensionality, from univalence;
\item the propositional truncation $\tr{A}$, as the image of $A\to 1$,
  with induction into propositions of any level and its computation rule
  as an identity;
\item set quotients of proposition-valued equivalence relations, as the
  image of $a \mapsto (x \mapsto R(x,a))$ into $A\to\mathsf{Prop}$
  (locally small by univalence); they are effective, and maps into sets
  that respect the relation descend to them;
\item univalent completions of small families.
\end{itemize}
The computation rules of the derived truncation and quotients are proved
identities. The evaluator has no new reduction rules.

\section{How the mathematics is represented}

\paragraph{Rings.} A commutative ring is a record with a carrier, the
operations, and the axioms as identities. Where needed the carrier is
assumed to be a set. A unit is an element with a chosen inverse, which is
a proposition when the carrier is a set.

\paragraph{Polynomials.} A polynomial is a list of coefficients, lowest
degree first. Two lists are equal as polynomials when all their
coefficients are equal, so trailing zeros do not matter. In the files on
$\PP$ the type $R[X]$ is the quotient of lists by this relation; in the
file on Lemma~2.2.4 we keep lists and the relation, which suffices.
Evaluation is Horner's rule, a primitive recursion. Laurent polynomials
are pairs of lists $(p,n)$, standing for $p(X) + X^{-1}n(X^{-1})$.

\paragraph{Recursion.} All recursive definitions and all proofs by
induction are written with the eliminators of the data types, as in
system~T but with dependent types. Several proofs use recursion at
higher type. For instance, the induction on the degree in
Lemma~2.2.4 returns, at each step, a function over all polynomials of
the smaller length.

\paragraph{Finitely presented algebras.} An $R$-algebra is a commutative
ring whose carrier is a set, with a ring homomorphism from $R$. A
presentation is a number $n$ and a list of expressions in the variables
$\mathsf{Fin}\,n$; expressions are an ordinary data type (variables,
constants of $R$, sums, products), and $\mathsf{Fin}$ is defined by
recursion ($\mathsf{Fin}\,0=\emptyset$, $\mathsf{Fin}(n+1)=\mathsf{Fin}\,n+1$),
so no indexed families are needed. The solutions of $P$ in an algebra $B$
are the assignments $V\to B$ on which every equation evaluates to $0$.
An algebra $A$ is \emph{presented} by $P$ when it has a solution $a$ of
$P$ such that for every algebra $B$ the map
$\mathrm{Hom}(A,B)\to\mathrm{Sol}_P(B)$, $\varphi\mapsto\varphi\circ a$, is
an equivalence; $A$ is \emph{finitely presented} when merely there is a
presentation of it, as in~\cite[Definition 1.3.2]{SAG}. Then
$\Spec A=\mathrm{Hom}(A,R)$ is equivalent to $\mathrm{Sol}_P(R)$, the
solutions of the equations in $R$. Every presentation has a presented
algebra: expressions modulo the relation of merely having a finite
derivation from the ring laws and the equations. A derivation is an
ordinary tree with a validity predicate defined by recursion. Two algebras
with the same presentation are isomorphic. The polynomial ring $A[X]$,
built as a quotient of lists, has the universal property
$\mathrm{Hom}(A[X],B)\simeq\mathrm{Hom}(A,B)\times B$, and is presented by
any presentation of $A$ with one more variable.

\paragraph{The axioms of \cite{SAG}.} Loc, SQC and Z-choice are not
part of the foundation. When a proof needs one of them, it is a
hypothesis of the segment containing the proof. SQC is stated once: for
every finitely presented algebra $A$, evaluation
$A\to(\Spec A\to R)$ is an equivalence. Some earlier files use
instances of SQC for particular algebras instead; this is said for each
result. Z-choice is used only in Theorem~7.3.6, where it is a hypothesis
for one affine $X$, together with weak quasi-coherence of the module.

\section{The results}

\subsection*{Proposition 2.2.2: the weak Nullstellensatz}

For a finitely presented $A$, $\Spec A=\emptyset$ if and only if $1=0$
in $A$. By SQC, an empty spectrum gives $A\simeq(\emptyset\to R)$, a
one-element ring; the converse uses $1\ne0$ in $R$.

For $A=R[X]/(f)$, presented by one variable and the equation $f$, the
algebraic content is a certificate: if $1=0$ in $A$ then merely $fg=1$
for some list $g$, with equality of coefficients. It is obtained by
reading a finite derivation of $1=0$ as a computation of a multiplier:
each rule (the ring laws, the equation $f=0$, symmetry, transitivity,
compatibility with $+$ and $\times$) transforms multipliers $r$ with
$x=y+fr$. No second quotient of $R[X]$ by an ideal is built, and SQC is
not used for this step. Combined with SQC: if $f$ has no root then merely
$fg=1$ for some $g$.

\subsection*{Lemma 2.2.4: $R$ is algebraically closed up to double negation}

For a polynomial $f$, it is not not the case that $f = 0$ or
\[
  f = \alpha\,(X-a_1)^{e_1}\cdots(X-a_n)^{e_n},
\]
with $\alpha$ a unit, $e_i\ge 1$ and the $a_i$ pairwise distinct.

The proof follows the paper: induction on a bound for the degree, inside
the double negation monad. The top coefficient is zero (drop it) or not.
In the second case it is a unit, and if the degree is at least $1$ there
is not not a root $a$. Then synthetic division by $X-a$ gives $f=(X-a)q$,
and $q$ has a smaller bound. Synthetic division is one primitive
recursion returning both the quotient and the remainder, and gives the
factor theorem with the bound on the degree. A last induction groups
equal roots, deciding $a_i=a_j$ under the double negation. This uses
that $X-a$ and $X-b$ commute, a coefficient computation.

\emph{Hypotheses:} Loc and SQC. The proof was first written with two
hypotheses, $1\ne0$ and the weak Nullstellensatz for $R[X]/(f)$ in the
form: if $f$ has no root then it is not not the case that $fg=1$ for
some list $g$. Double negation suffices because the proof only uses it
inside double negations, and the mere inverse given by
Proposition~2.2.2 implies it. With Proposition~2.2.2 the lemma follows
from Loc and SQC alone. From the Nullstellensatz we prove
Proposition~2.2.3(a) up to double negation (a constant $x\ne 0$ has no
root, hence is not not invertible). We also prove that a polynomial of degree $\ge 1$ with a
unit top coefficient is not invertible: the top coefficient of $fg$ is
the product of the top coefficients. Replacement is not needed here.
Nine computations on $\mathbb{Z}/2$ check that the definitions compute,
for instance $1+X^2=(X+1)^2$.

\subsection*{Section 3.1: affine types, localisation and affine opens}

An affine witness for $X$ consists of an $R$-algebra $A$, mere finite
presentability of $A$, and an equivalence $X\simeq\Spec A$;
affineness is its propositional truncation (Definition~3.1.1).
The pointwise ring $R^X$ is constructed as an actual $R$-algebra.
Under SQC, evaluation at points gives the canonical equivalence
$X\simeq\Spec(R^X)$ (Lemma~3.1.2).
For Proposition~3.1.3, both representing algebras are compared with
$R^X$. A structure identity lemma, proved from univalence, retains
the $X$-indexed points as well as the ring operations and constants.
It therefore identifies the entire affine witnesses, including their
spectrum identifications. Their type is a proposition. In particular,
$R^X$ is finitely presented when $X$ is affine.

\paragraph{The algebraic constructions.}
Adjoining a variable $y$ with equation $fy-1$ constructs $A[f^{-1}]$,
with finite presentability and its canonical universal property for
every target $R$-algebra $B$:
\[
 \mathrm{Hom}(A[f^{-1}],B)
 \simeq \sum_{h:\mathrm{Hom}(A,B)}\mathrm{inv}(h(f)).
\]
At $B=R$ this gives $D(f)=\Spec A[f^{-1}]$ (Proposition~3.1.6).
Appending finitely many equations constructs $A/(g_1,\ldots,g_n)$
with the analogous universal property and spectrum the simultaneous
zero locus. Elements of a presented algebra are merely expressions in
its generators. Those expressions, and finite presentations themselves,
are eliminated only into propositions: no presentation is chosen from
mere finite presentability.

\paragraph{Opens and certificates.}
Uniform SQC and $1\ne0$, applied to $R/(a)$, imply
$a\ne0\Rightarrow\mathrm{inv}(a)$; the converse is elementary.
This gives Example~3.1.4 and the equivalence between the two descriptions
of affine-open subtypes in Definition~3.1.5.
Every element of $A[f^{-1}]$ is merely a fraction $h/f^k$.
The open it defines is the restriction of $D(hf)$, which is contained
in $D(f)$. Finite choice of fractions proves Lemma~3.1.7, stated also
as affine-openness of the direct image under
$\Spec A[f^{-1}]\to\Spec A$. No geometric axiom is needed for this
transitivity proof.

For Lemma~3.1.8, collapse of $A[f^{-1}]$ gives a polynomial certificate
in $A[Y]/(fY-1)$. Comparing coefficients extracts an exponent with
$f^k=0$. With SQC and $1\ne0$, this proves
$D(f)=\emptyset$ if and only if $f$ is nilpotent.
For Lemma~3.1.9, quotient by the $g_i$, apply Lemma~3.1.8 to the
image of $f$, and interpret the quotient's finite derivation trees as
ideal certificates. The result is the equivalence
\[
 D(f)\subseteq\bigcup_i D(g_i)
 \quad\Longleftrightarrow\quad
 \left\lVert\sum_{k:\mathbb N}\sum_{c:\mathsf{Fin}(n)\to A}
       f^k=\sum_i g_i c_i\right\rVert .
\]
The reverse direction uses Loc to find an invertible summand.
Thus the combined formal statement has explicit Loc and uniform SQC
hypotheses. The final corollary is also proved: the $D(g_i)$ cover
$\Spec A$ exactly when the $g_i$ generate the unit ideal.

\subsection*{Sections 3.2 and 3.3.1--3.3.2: tensor products, pullbacks, eventually vanishing sequences}

\emph{Lemma 3.2.1.} The tensor product $A\otimes B$ is the coproduct of
$R$-algebras: maps $A\to T$, $B\to T$ such that $h\mapsto(h\circ
\iota_A,h\circ\iota_B)$ is an equivalence $\mathrm{Hom}(T,D)\simeq
\mathrm{Hom}(A,D)\times\mathrm{Hom}(B,D)$ for every $D$. With $D=R$,
$\Spec T=\Spec A\times\Spec B$. For $A$, $B$ presented by $p$ and $q$,
the algebra presented by the concatenation of $p$ and $q$ (variables of
both, equations of both, renamed) is a coproduct; so finitely presented
algebras have a finitely presented tensor product. No SQC is used.

\emph{Lemma 3.2.2.} For $f:C\to A$ and $g:C\to B$, add to the
concatenated presentation one equation $e_v-e'_v$ for each generator $c_v$
of $C$, where $e_v$, $e'_v$ are expressions in the generators for
$f(c_v)$ and $g(c_v)$. The algebra presented is a pushout, and its
spectrum is the pullback $\Spec A\times_{\Spec C}\Spec B$. The
expressions exist merely: every element of a presented algebra is merely
the value of an expression, and finitely many can be chosen at once. By
SQC every map $\Spec A\to\Spec C$ comes from a unique algebra map
$C\to A$, so the pullback of any two maps of affine schemes is affine.

\emph{Lemma 3.3.1.} For a tensor product $T$ of $A$ and $B$, with $T$ and
$B$ finitely presented, $c\mapsto(\varphi\mapsto(\varphi\otimes B)(c))$
is an equivalence $T\simeq(\Spec A\to B)$. The proof is the paper's
chain: SQC for $T$, Lemma~3.2.1, currying, SQC for $B$, using
$\psi\circ(\varphi\otimes B)=\varphi\otimes\psi$.

\emph{Lemma 3.3.2.} $A[X]$, the polynomial ring of the polynomial
library, is a tensor product of $A$ and $R[X]$. An eventually vanishing
sequence is the sequence of coefficients of a unique element of $R[X]$;
uniqueness lets us build it from a mere bound, without choice. Lemma~3.3.1
with $B=R[X]$ turns the family into one polynomial $p$ over $A$, and
$x\otimes R[X]$ applies $x$ to its coefficients. The length of a list
representing $p$ is the common bound.

\subsection*{Theorem 3.3.3: functions $\Spec A\to\mathbb N$ are bounded}

For a finitely presented $A$ and $f:\Spec A\to\mathbb N$, merely there is
$N$ with $f(x)\le N$ for all $x$. The paper's proof goes through tensor
products (Lemma~3.3.1) and eventually vanishing sequences. A shorter route
uses only $A[T]$:
\begin{itemize}
\item the function $H(x,t)=t^{f(x)}$ on $\Spec A\times R$ is a function on
  $\Spec(A[T])$, since a homomorphism $A[T]\to R$ is a point of $\Spec A$
  and the image of $T$;
\item by SQC for $A[T]$, $H$ is merely the evaluation of one polynomial
  $p$ over $A$;
\item at a point $x$, the polynomial obtained by applying $x$ to the
  coefficients of $p$ has the same values as $T^{f(x)}$, hence the same
  coefficients (SQC for $R[T]$);
\item if $f(x)$ is at least the length of $p$, its coefficient at
  $f(x)$ is $0$, while that of $T^{f(x)}$ is $1$; this contradicts
  $1\ne0$.
\end{itemize}
The bound is the length of $p$. It is a mere bound because $p$ is merely
a list. Order on $\mathbb N$ is decidable, so no excluded middle is used.
There are two versions: one with $1\neq0$ and the two halves of SQC used
(surjectivity for $A[T]$, injectivity for $R[T]$) as hypotheses, and one
from $1\ne0$ and uniform SQC, where $A[T]$ and $R[T]$ are shown finitely
presented.

The paper's proof is also formalised: it applies Lemma~3.3.2 to the
sequences equal to $1$ below $f(x)$ and $0$ from $f(x)$ on, and uses
$1\ne0$. It has the same hypotheses as the uniform version of the
shorter proof: uniform SQC, $1\neq 0$ and finite presentation of $A$.

\subsection*{Proposition 3.3.4: upwards closed families have a common element}

The axioms of Section~2.1 are now stated
uniformly: Loc, SQC and Z-choice, with an independent universe for the
family in Z-choice. For an upwards closed proposition-valued family
$P(x,n)$ on $\Spec A\times\mathbb N$, Z-choice turns pointwise mere
inhabitants into finitely many local sections on $D(f_i)$, where the
$f_i$ form a unimodular family. Proposition~3.1.6 and Theorem~3.3.3
merely bound each section. Finite choice combines the bounds; their
sum bounds them all. The covering property and upward closure give
\[
 \left(\prod_x\left\lVert\sum_n P(x,n)\right\rVert\right)
 \longrightarrow
 \left\lVert\sum_N\prod_x P(x,N)\right\rVert .
\]
There is no choice over the points of $\Spec A$. The implementation
uses the shorter proof of Theorem~3.3.3.

The combined entry point \kw{affine/Section3.seg} checks both
developments together: 2646 declarations, no failures, in $0.83$ seconds.
Integration lemmas turn the tensor and arbitrary-map pullback
constructions into the intrinsic predicate of affineness.

\subsection*{Section 4.1: closed subtypes and gluing ideals}

A closed proposition is merely a finite system of equations over $R$:
an equivalence with $\bigwedge_i x_i=0$ (Definition~4.1.1).
Ideals are proposition-valued predicates with the usual closure laws;
finite generation is mere existence of a finite family generating the
predicate. Membership in a generated ideal is the truncation of an
explicit linear-combination certificate.

\emph{Proposition 4.1.2} (SQC). Finitely generated ideals of $R$ are
order-reversingly equivalent to closed propositions. The forward map
sends $I$ to ``every member of $I$ is zero'', proved equivalent to
$I=(0)$. The inverse sends a closed proposition $Q$ to the ideal
\[
                   \{a:R\mid Q\to a=0\}.
\]
This constructs the predicate without choosing an equation witness from
a truncation. Both orders are partial orders on sets, and the equivalence
reverses the orders.

\emph{Lemma 4.1.3} (SQC). Inclusion
$V(f_1,\ldots,f_n)\subseteq V(g_1,\ldots,g_m)$ is equivalent to
$(g_1,\ldots,g_m)\subseteq(f_1,\ldots,f_n)$ as actual ideals.
Indeed, SQC on $A/(f_1,\ldots,f_n)$ detects that the images of the $g_j$
vanish; quotient certificates then recover their ideal membership.
This proof does not pass to radicals. A regression using dual numbers
checks the distinction: $\epsilon^2=0$, but $\epsilon$ is not in the
zero ideal; under SQC, $V(0)$ is not contained in $V(\epsilon)$.

\paragraph{Lemma 1.3.8: the algebraic gluing step.}
First, from the canonical localisation universal property, we prove
\[
 a/1=0\text{ in }A_f
 \quad\Longleftrightarrow\quad
 \left\lVert\sum_{k:\mathbb N} f^k a=0\right\rVert ,
\]
and its ideal-membership version
\[
 a/1\in(g_j/1)
 \quad\Longleftrightarrow\quad
 \left\lVert\sum_{k:\mathbb N} f^k a\in(g_j)\right\rVert .
\]
The first uses a polynomial certificate in $A[Y]/(fY-1)$ and coefficient
comparison; the second combines a quotient with this localisation model.
Neither needs a geometric axiom.

For compatible finitely generated ideals in the $A_{f_i}$, replace local
generators by numerators. Clear denominators on each overlap, combine
the finitely many exponents, and concatenate the numerators multiplied
by these powers of $f_i$. The resulting ideal restricts to each
prescribed ideal. Ideal extension is proved independent of generating
witnesses. The final statement uses literal equalities of extended
ideals. Iterated localisations are shown to satisfy the universal
property of $A_{f_if_j}$. The ring $A$ is arbitrary, and the theorem
does not require that the $D(f_i)$ cover; there is no uniqueness claim
in that generality. Localisation witnesses carry finite presentability
over $A$, which is derived from their universal property rather than
assumed separately.

\paragraph{Lemma 4.1.4.}
For a pointwise closed subtype $C$ of $\Spec A$, Z-choice gives local
finite equation witnesses on a finite unimodular cover. Boundedness of
their lengths allows padding with zero equations; SQC turns the
coordinates into elements of the local algebras. On overlaps their
zero loci agree, so Lemma~4.1.3 gives agreement of their ideals.
Lemma~1.3.8 glues these to a finitely generated ideal of $A$.
The cover then identifies the global zero locus with $C$.
The checked result retains a quotient $Q=A/(g_1,\ldots,g_n)$ with its
universal property and an equivalence
\[
                     \sum_{x:\Spec A} C(x)\simeq\Spec Q
\]
\emph{over} $\Spec A$. A separate theorem states that a closed subtype
of any affine type is affine. The hypotheses are Loc, uniform SQC and
Z-choice. There is no choice over the points of $\Spec A$.

The combined entry point \kw{closed/Section4.seg} also checks the open
subtype development below, with the same definition of closed propositions:
3860 declarations, no failures, in $1.06$ seconds.

\subsection*{Section 4.2: open subtypes}

A proposition is open if it is merely equivalent to one of finitely many
elements of $R$ being invertible (Definition~4.2.1); this has the same
form as closed propositions, with invertibility in place of vanishing.

\emph{Proposition 2.2.3(b),(c).} A vector is nonzero if and only if one of
its entries is invertible: $R/(x_1,\ldots,x_n)$ has no points, so it
collapses by Proposition~2.2.2, its kernel certificate gives
$1=\sum x_ic_i$, and Loc finds an invertible summand. An element is
nilpotent if and only if it is not not zero, by Lemma~3.1.8 for $A=R$.

\emph{Propositions 4.2.2--4.2.3, Lemma 4.2.4.} Open propositions contain
false and true, are closed under binary unions (append the lists) and
intersections (all products), and pull back along maps; no axiom is used.
By 2.2.3(b) they are the negations of closed propositions, hence
double-negation stable (Loc, SQC).

\emph{Lemmas 4.2.5--4.2.6.} On an affine $X$, $x\ne y$ is open: a
presentation embeds $\Spec A$ in $R^n$, and $x\ne y$ says that the vector
of coordinate differences is nonzero. So $x$ specialises to $y$ if and
only if $\neg\neg(x=y)$.

\emph{Theorem 4.2.7} (Loc, SQC, Z-choice). An open subtype $S$ of
$\Spec A$ is affine-open. As in the paper: Z-choice gives open witnesses
on a cover $D(f_i)$; their lengths are bounded on
$D(f_i)=\Spec A[f_i^{-1}]$ by Theorem~3.3.3; padding with zeros and SQC
for $A[f_i^{-1}]$ give elements $g_{ij}$ with $S\cap D(f_i)=D(g_{ij})_j$;
Lemma~3.1.7 brings these down to $\Spec A$, and the lists are
concatenated. The last step is a gluing lemma used again below: a
subtype covered by finitely many $D(f_i)$ and affine-open on each
$\Spec A[f_i^{-1}]$ is affine-open.

\emph{Lemmas 4.2.8--4.2.9.} In $\Spec A$, an open subtype of an open
subtype is open: by 4.2.7 the outer one is $D(f_1,\ldots,f_n)$, on each
$\Spec A[f_j^{-1}]$ the inner one is open, hence affine-open, and gluing
applies. The case $A=R$ says that open propositions are closed under
dependent sums (4.2.9(a)), and applying it pointwise gives transitivity
of open subtypes for every type (4.2.9(b)), schemes included, without a
definition of schemes.

\emph{Lemma 4.2.11.} For $P$ closed and $Q$ open,
$(P\to Q)\Leftrightarrow\neg P\vee Q$. The paper states no axioms for it,
but its proof uses Proposition~2.2.3(b), so Loc and SQC are hypotheses
here.

\subsection*{Section 5: schemes}

\paragraph{Small covers and local choice.}
A chart of $X:\U_l$ consists of a presentation $p$, an embedding
$\operatorname{Sol}_p(R)\to X$, and openness of its image. A cover
consists of finitely many charts and a proof that every point merely
belongs to one of their images. Both $\operatorname{Cover}(X)$ and
\[
              \operatorname{Scheme}(X)=\tr{\operatorname{Cover}(X)}
\]
belong to $\U_l$. No setness hypothesis is imposed on $X$: it follows
from the chart embeddings and the setness of solution types.
Finite choice proves equivalence with the paper's definition by affine
open subtypes. The key reduction is that the large assertion that $T$
is affine is equivalent to the small proposition
$\tr{\sum_p T\simeq\operatorname{Sol}_p(R)}$.

For Proposition~5.1.3, apply Z-choice on the presented algebra of each
chart to a family $B:X\to\U_l$ with pointwise mere inhabitants.
This gives finite unimodular lists of functions and sections on their
principal opens. Localisation identifies each principal open with an
affine type, which can again be represented by a small presentation.
Finite choice and concatenation yield merely a finite affine open cover
of $X$ with sections of $B$. The proof retains the standard-open
refinement data, and also provides sections on the chart images.
All Z-choice applications are at level $l$; no higher-universe choice
or choice over all points of $X$ is added.

A related lemma flattens a finite open cover by schemes to an affine
open cover. It uses finite choice of the covers, concatenation, and
open transitivity. It does not separately invoke Theorem~4.2.7.

\paragraph{Specialisation and gluing.}
For Lemma~5.2.1, take an affine chart containing $x$.
If $x$ specialises to $y$, the chart also contains $y$.
Open transitivity extends opens of the chart to $X$, so the affine
specialisation theorem gives $\neg\neg(x=y)$.
The reverse implication holds for any type.

For Proposition~5.3.1, given embeddings $f:V\to X$ and $g:V\to Y$
into sets, quotient $X+Y$ by equality on each summand and by
\[
 \operatorname{Both}(x,y)=
       \sum_{v:V}(f(v)=x)\mathbin{\times}(g(v)=y)
\]
between the summands. The embedding assumptions make this an equivalence
relation with proposition-valued fibres. Effectiveness of the quotient
shows that its two inclusions are embeddings and identifies their
intersection. Quotient descent proves that restriction
\[
 (P\to B)\ \longrightarrow\
 \sum_{h:X\to B}\sum_{k:Y\to B}\prod_{v:V} h(f(v))=k(g(v))
\]
is an equivalence for every set $B:\U_l$. The universal property
against arbitrary types, or sets at a separate larger level, is not
claimed.

When $f$ and $g$ have open images, the two inclusions into $P$ do too:
this is checked on quotient representatives by proposition-valued
induction. Composing the affine charts of $X$ and $Y$ with these open
embeddings and concatenating them gives a cover of $P$.
The quotient and its set-target universal property use no geometric
axiom; openness uses the proved transitivity theorem.

\paragraph{Subschemes and equality types.}
For Propositions~5.4.2--5.4.3, restrict a scheme's affine charts to the
given subtype. For a closed subtype, each restriction is affine by
Lemma~4.1.4. For an open subtype of an affine type, Theorem~4.2.7
provides a finite cover by principal opens, each affine by localisation.
The finite-cover flattening lemma gives schemehood in the general
open case. The definition of subscheme explicitly includes
proposition-valuedness of the family.

Lemma~5.5.1 uses no geometric axiom: equality in an affine type is a
finite system of coordinate equations. Every closed proposition is
itself affine, using a presentation with zero variables and those
constant equations.
For Proposition~5.5.2, choose an affine open $U$ containing $x$, with
witness $u_x$. The open proposition $U(y)$ is a scheme, and
\[
             \sum_{u:U(y)} (x,u_x)=(y,u)
\]
is its closed subtype, where the equality is in the affine chart.
It is therefore a scheme and is equivalent to $x=y$.
The argument allows $U(y)$ to be empty.

\paragraph{Dependent sums and families (Section 5.6).}
Theorem~5.6.1 states that, for a scheme $X$ and a family of schemes
$Y:X\to\U_l$, the total space $\sum_{x:X}Y(x)$ is a scheme.
The proof uses two distinct padding constructions.
First suppose $X$ is affine and each fibre has a chosen presentation.
Theorem~3.3.3 merely bounds their numbers of variables by some $N$.
Each solution set is equivalent to a closed subtype of $R^N$:
restrict to the original variables, impose the original equations,
and require all added coordinates to be zero. The numbers of equations
may vary without any uniform bound. Consequently the total space is a
closed subtype of the affine type $X\times R^N$, hence affine by 4.1.4.
Local choice of presentation witnesses and finite-cover flattening then
show that an arbitrary family of affine fibres over a scheme has a
scheme total space. This last assertion does not claim affineness.

For general scheme fibres, local choice gives actual finite covers of
the fibres on affine charts of the base. Theorem~3.3.3 bounds the numbers
of charts locally. Pad each cover to that bound using empty affine
charts, presented by the equation $1=0$. For each fixed chart index,
the family has affine fibres, so the preceding result applies.
The resulting finite open cover proves schemehood of the total space.
All choices are eliminated into proposition-valued goals and all
Z-choice applications remain at level $l$.

Univalence gives the equivalence between maps into $X$ and $X$-indexed
families: a map is sent to its actual fibres, and a family to its total
space with its projection. Closure under dependent sums and under
scheme-valued equality restricts this to the equivalence of schemes in
Corollary~5.6.2. For proposition-valued families this gives 5.6.3:
a subtype is a subscheme exactly when every value is a scheme.
Finally, the pullback
\[
             \sum_{x:X}\sum_{y:Y} f(x)=g(y)
\]
is a scheme by two applications of dependent-sum closure and 5.5.2,
proving 5.6.4.

\paragraph{Scope and checks.}
The formal geometric proofs take Loc, SQC and Z-choice.
In particular, the present proofs of 5.1.3 and 5.2.1 use open
transitivity and do not recover the paper's weaker hypothesis lists
(Z-choice alone and Loc/SQC, respectively).
The combined entry point \kw{schemes/Section5.seg} checks
5194 declarations with no failures in $1.30$ seconds.
Regressions include empty covers, empty subtypes, and the empty
equality type between zero and one on the affine line, zero-constrained
padding, empty pullbacks, and classification of a map with two-point
fibres. All 20 expected negative tests are rejected.

\subsection*{Example 6.1.9: squaring on the projective line}

$\PP$ is the quotient of two copies of $R$ (the charts $[a:1]$ and
$[1:b]$) by the relation $ab=1$. It is shown equivalent to the quotient
of unimodular pairs by multiplication with units. Squaring descends to
$s:\PP\to\PP$, with $s([x:y])=[x^2:y^2]$. The two fibres are computed:
\[
  s^{-1}([0:1]) \simeq \{a:R\mid a^2=0\},
  \qquad
  s^{-1}([1:1]) \simeq \mathsf{Bool}\quad\text{if } 2\ne 0 .
\]
The first is the case $t=0$ of a chart computation valid for every
$t$, $s^{-1}([t:1])\simeq\{a\mid a^2=t\}$, which only needs a ring. The
second uses Loc, $1\neq 0$, and that $2$ is invertible. That $2\neq 0$
implies $2$ invertible uses SQC for the algebras $R/(x,y)$, given as a
quotient of $R$, with evaluation into functions on $\{x=0\wedge y=0\}$.
These are the algebras used in the proof of Proposition~2.2.3.

The coordinate algebra of a fibre is treated in the setting of
presentations. Let $K$ be the algebra presented by one variable $T$ and
the equation $T^2=a$. Then $(b,c)\mapsto b+cT$ is an equivalence
$R\times R\simeq K$ compatible with the module operations, so $K$ is free
of rank two with basis $1,T$. With SQC, the functions on
$\{x:R\mid x^2=a\}$ are exactly the $x\mapsto b+cx$. For $a=0$ this is
first-order expansion on the infinitesimal disk: every $f$ satisfies
$f(d)=f(0)+f'\,d$ for a unique slope $f'$. This uses neither Loc nor the
invertibility of $2$. The connection with the fibres of $s$ computed
above is not made formally.

\subsection*{Lemma 6.2.1: functions $\PP\to R$ are constant}

The map $R\to(\PP\to R)$ sending $c$ to the constant function is an
equivalence, with inverse evaluation at $[0:1]$. The hypotheses are SQC
for $R[X]$ (evaluation $R[X]\to(R\to R)$ is an equivalence) and SQC for
$R[X,X^{-1}]$ (evaluation into functions on $\{(x,y)\mid xy=1\}$). The two
restrictions of $f$ to the charts are merely polynomials $p(X)$ and
$q(X^{-1})$. They agree on the overlap, so by injectivity for the
Laurent polynomials all positive coefficients vanish and the constants
agree. The coefficient computation does not use SQC. Constancy then
follows by induction over the quotient $\PP$.

\subsection*{Proposition 4.6 of \cite{Proj}: line bundles on $\PP$}

$\mathsf{Line}$ is the small type of $R^\times$-torsors and $\PP$ the
quotient of two charts above. The map
$\mathsf{Line}\times\mathbb Z\to(\PP\to\mathsf{Line})$,
$(l_0,d)\mapsto(x\mapsto l_0\otimes\mathcal O(d)(x))$, is an equivalence,
from Loc, SQC and Z-choice. Its inverse sends $L$ to $L([0:1])$ and the
degree of $L$; both inverse equations are checked.

\paragraph{Horrocks (Lemma 2.1 of \cite{Proj}).} A line bundle on
$\Spec A$ is given by a rank-one projector $E$ over $A[X]$. Let $A$ be
any commutative ring and $f$ monic. If $E$ factors after inverting $f$,
that is $x_iy_j=f^kE_{ij}$ and $\sum_iy_ix_i=f^k$, then $E=ba$ with
$\sum_ia_ib_i=1$. No SAG axiom is used. The proof follows Appendix~1
of~\cite{Proj}, with explicit certificates throughout:
\begin{itemize}
\item \emph{Formal gcd.} A node is a pair of finite lists $(I,U)$ of
  elements of $A$; it inverts the monoid $S(I;U)=M(U)+(I)$, which puts
  $(I)$ in the Jacobson radical. Starting from $f^k$, the row ideal
  $(y_1,\dots,y_n)$ always contains a monic polynomial. The remainders of
  the $y_i$ are scanned from the top coefficient down. A coefficient $a/s$
  is tested by branching on $a$: either $a$ joins $I$, or $a$ joins $U$
  and becomes a unit. A unit coefficient gives, by monic lifting, a monic
  member of smaller degree. The recursion is on the degree bound, the
  position and the number of coefficients; no degree or unit is decided.
\item \emph{Leaves.} At each leaf the monic member generates the row
  ideal modulo $(I)$. By the saturation calculation and a determinant-free
  Nakayama elimination in $A_S[X]/(h)$, it generates the row ideal itself.
\item \emph{Comaximality.} For any choice of one $s_\lambda\in S_\lambda$
  per leaf, the $s_\lambda$ generate the unit ideal: two children combine
  through $w=s^kv+(-b)^kj$, and the root monoid is $\{1\}$.
\item \emph{Gluing.} Clearing denominators at leaf $i$ gives
  $d_i,\tau_i\in S_i$ and, in $A[X]$, $p_i=\sum_jy_jz_{ij}$ with
  $p_i=d_ih_i$ over $A_{S_i}$ and $\tau_iy_j=p_ib_{ij}$. Over the product
  monoid $S_iS_j$ both monic generators generate the row ideal, so they
  are equal (also when their degrees differ, since then $1=0$); hence
  $u_{ij}v_{ij}(d_jp_i-d_ip_j)=0$ with $u_{ij}\in S_i$, $v_{ij}\in S_j$.
  With $r_i=\prod_ku_{ik}v_{ki}$ the numerators $H_i=r_ip_i$ satisfy
  $s_iH_j=s_jH_i$ for $s_i=d_ir_i$, and $t_iy_j=H_ib_{ij}$ for
  $t_i=r_i\tau_i$. Two partitions of unity, for $(s_i)$ and $(t_i)$, give
  $g\in(y)$ with $y_j=ga_j$. The global $g$ need not be monic.
\item \emph{Descent.} From $g=\sum_jy_jc_j$ and $y_j=ga_j$, put
  $b_i=\sum_jE_{ij}c_j$; cancelling the regular element $f^k$ gives
  $E_{ij}=b_ia_j$ and $\sum_ia_ib_i=1$.
\end{itemize}

\paragraph{Chart triviality (Corollary 2.3 of \cite{Proj}).} A line bundle
on an affine algebra gives a rank-one projector whose fixed vectors are
the global sections. A frame of $L$ on $D(f)$, extended by Lemma~7.1.8,
gives exactly the certificate above, so Horrocks makes $L$ trivial on
$\mathbb A^1$ when it is trivial on a monic open. Z-choice gives a
polynomial neighbourhood of $0$ with a frame; reversing and normalising a
polynomial with invertible constant term gives a monic open in the other
chart. Hence the restriction of $L:\PP\to\mathsf{Line}$ to each chart is
trivial.

\paragraph{Laurent algebra (Lemmas 4.1 and 4.4 of \cite{Proj}).} On the
overlap, $L$ is then given by a function from units to units. By SQC
every function on the units is merely a Laurent polynomial. Lemma~4.1
is proved for $R$: an invertible Laurent polynomial merely has one unit
coefficient, and every other coefficient is not nonzero. Coefficients are
decided only under double negation. Lemma~4.4 factors an invertible Laurent polynomial as
$uX^d(1+a)(1+b)$ with $a\in XR[X]$, $b\in X^{-1}R[X^{-1}]$ nilpotent; the
nilpotent part is removed one power of a single nilpotent at a time.
The degree is unique, since $x^Na(x)=x^Mb(1/x)$ with units $a,b$ forces
$N=M$.

\paragraph{Scope and checks.} The Horrocks package checks in 2.79
seconds (4982 declarations), the Laurent algebra in 1.17 seconds, and
the whole of Proposition~4.6 in 4.84 seconds (9773 declarations, no
failures), also from an isolated copy of its sources. A false final proof
is rejected.

\subsection*{Section 7 of \cite{Proj}: line bundles on $\mathbb P^n$}

\paragraph{The two definitions.} $\mathbb P^n$ is defined as the vectors
of $R^{n+1}$ with an invertible coordinate, modulo multiplication by
units; the relation is a proposition and is equivalent to
proportionality. It is shown equivalent to the type of pairs of
$l:\mathsf{Line}$ and a nonzero equivariant map from $l$ to $R^{n+1}$.
With Loc and SQC, a vector has an invertible coordinate if and only if it
is nonzero. These results hold for any commutative ring.

\paragraph{Lemma 1.6 and Proposition 1.7.} For $k\ge2$, every function
$R^k\setminus\{0\}\to R$ extends uniquely to $R^k$. This is proved as a
Hartogs lemma for a finitely presented algebra $A$ with regular elements
$s_0,s_1,\dots$ such that $s_1^bg=s_0^ah$ implies $s_0^a\mid g$: functions
on $D(s_0)\cup D(s_1)\cup\cdots$ are elements of $A$. It follows from SQC
for $A$ and its localisations. Proposition~1.7, every function
$\mathbb P^n\to R$ is constant, follows as in~\cite{Proj}.

\paragraph{Lemma 7.1.} For $n\ge2$ and $p\ne q$, functions
$\mathbb P^n\setminus\{p,q\}\to R$ are constant. A linear automorphism,
built from two transpositions and Cramer's rule, moves $p,q$ to
$[e_Y],[e_Z]$. Then $\mathbb P^n\setminus\{[e_Y],[e_Z]\}$ lifts to the
open subset $D(X)\cup D(YZ)\cup\bigcup_jD(x_j)$ of $R^{n+1}$, $x_j$ the
other coordinates. This family satisfies the hypotheses of the Hartogs
lemma, so no gluing of chart extensions is needed. For
$\mathbb P^n\setminus\{p\}$, a point $x\ne[e_Y]$ is apart from $[e_Z]$
or from $[e_X]$, and the two cases share the point $[e_X+e_Z]$. The
statements for $\mathbb Z$ follow from those for $R$, since $\mathbb Z$ has
decidable equality.

\paragraph{$\mathcal O(d)$ and Lemma 7.2.} The fibre of $\mathcal O(d)$
at $p$ is the torsor of unit-valued functions $v$ on the representatives
of $p$ with $v(ux)=u^dv(x)$; so $\mathcal O(1)$ is dual to the tautological
line. For $[x]\ne[y]$ the line $\PP\to\mathbb P^n$ sends $[a:1]$ to
$[ax+y]$ and $[1:b]$ to $[x+by]$. The representatives of a point of $\PP$
in the two charts give representatives of its image with the same
transitions, so $\mathcal O(d)$ restricts to the $\mathcal O(d)$ of
Proposition~4.6. With $\mathrm{Gr}(1,M)$ parametrised by $\PP$ through a
basis $x,y$ of $M$, this is Lemma~7.2.

\paragraph{Theorem 7.3.} For $n\ge2$ the map
$\mathsf{Line}\times\mathbb Z\to(\mathbb P^n\to\mathsf{Line})$,
$(l_0,d)\mapsto(x\mapsto l_0\otimes\mathcal O(d)(x))$, is an equivalence.
Its inverse sends $L$ to $L([e_Z])$ and the degree of $L$ on the line
through $[e_Y]$ and $[e_Z]$. The main step is that two line bundles
$L_1,L_2$ on $\mathbb P^n$ that agree on one line are merely equal:
\begin{itemize}
\item For a point $p$, that $L_1$ and $L_2$ agree on the lines through
  $p$ and $x$ is a proposition of $x\ne p$. It is decidable, by comparing
  degrees on $\PP$, so by Lemma~7.1 it holds for all $x$ once it holds for
  one.
\item Since functions $\PP\to R^\times$ are constant, an identification
  $L_1(p)=L_2(p)$ extends uniquely along each line through $p$. This
  identifies $L_1$ and $L_2$ on $\mathbb P^n\setminus\{p\}$, independently
  of the representatives chosen for the line.
\item Doing the same at a second point $q$, the two identifications
  differ on $\mathbb P^n\setminus\{p,q\}$ by a function to $R^\times$,
  constant by Lemma~7.1. After correcting by it they glue.
\end{itemize}
This avoids $L\otimes\mathcal O(-d)$, so associativity and inverses for
the tensor product of lines are not needed. An equality between merely
equal bundles is determined by its value at one point, since functions
$\mathbb P^n\to R^\times$ are constant.

\paragraph{Corollary 5.6: $\mathrm{Aut}(\mathbb P^n)=\mathrm{PGL}_{n+1}(R)$.}
For $n\ge2$, the action of $\mathrm{PGL}_{n+1}(R)$, the linear automorphisms of
$R^{n+1}$ modulo units, on $\mathbb P^n$ is an equivalence onto the maps
$\mathbb P^n\to\mathbb P^n$ that have an inverse, and it respects
composition. The proof does not use Corollary~5.5 or tensor powers of
lines.
\begin{itemize}
\item A map $\varphi:\mathbb P^n\to\mathbb P^m$ merely lifts to
  $F:R^{n+1}\setminus\{0\}\to R^{m+1}\setminus\{0\}$ with
  $F(ux)=u^dF(x)$: in the first definition, $\varphi$ gives a line
  bundle with a nonzero equivariant vector; by Theorem~7.3 the bundle is
  merely $\mathcal O(d)$, and a representative $x$ gives a frame of
  $\mathcal O(d)$ at $[x]$.
\item By Lemma~1.6 a function of negative weight vanishes, so $d\ge0$; and
  a function that is a unit on $R^{k}\setminus\{0\}$, $k\ge2$, extends to a
  unit (extend it and its inverse), so it cannot have positive weight.
\item For $\varphi$ with inverse $\psi$ and lifts $F,G$ of weights $a,b$,
  $G(F(x))=\lambda(x)x$ with $\lambda(ux)u=(u^a)^b\lambda(x)$. By the
  previous point $a=b=1$, and $\lambda$ is constant by Proposition~1.7.
\item A function of weight $1$ is additive: $f(x+ty)=f(x)+ta_1(t)$, and
  the chart intersection lemma of Lemma~6.2.1 makes $a_1$ constant,
  equal to $f(y)$. So $F$ extends to a linear map, with $GF$ and $FG$
  scalars.
\item Injectivity: if $[Ax]=[Bx]$ for all $x$, the ratio is a function on
  $\mathbb P^n$, hence a constant unit.
\end{itemize}

\paragraph{Comparison with \cite{Proj}.} The formal proofs supply some
steps that the text leaves implicit.
\begin{itemize}
\item Corollary~5.6 is stated without proof after Corollary~5.5. The
  missing steps are: the degrees $d,e$ of an automorphism and of its
  inverse satisfy $de=1$; a map of degree $1$ is linear; it is invertible
  up to a scalar; and two matrices inducing the same map differ by a unit.
\item In Corollary~5.5 the degree is not negative. This needs an argument:
  a function of negative weight on $R^{n+1}\setminus\{0\}$ vanishes.
\item In Lemma~7.1, $U_1$ should be $\{x_1\ne0\}$. That one may assume
  $p\in U_0\setminus U_1$ and $q\in U_1\setminus U_0$ needs a change of
  coordinates, and $n>1$ is where Lemma~1.6 is used. Parts (i) and (iii)
  need a third point, or the observation that every $x\ne p$ is apart
  from one of two further points.
\item In Lemma~7.2 the map should be $L\mapsto l_0\otimes L^{\otimes d}$,
  and for $d<0$, $\mathrm{Hom}(L^{\otimes(-d)},R)$. In Theorem~7.3, $n>1$
  is needed, and independence of the degree needs both that it is
  symmetric in $p,q$ and that every $x$ is apart from $p$ or from $q$.
\end{itemize}

\paragraph{Scope and checks.} The directory \kw{pn} checks 12188
declarations with no failures in 6.6 seconds, including all of
\kw{picard}. Each of its files checks in at most 0.25 seconds. Other
results of~\cite{Proj} are not formalised.

\subsection*{Section 7.1: weakly quasi-coherent modules on schemes}

\paragraph{Modules.} An $R$-module is the module record of the torsor
files, over the same carrier and operations as $R$. Sections~2--5 and
the torsor files use separate copies of $\mathbb N$ and of finite types;
explicit translations between them make powers and finite sums agree.
For $X:\U_l$ and modules $M_x$, $M(X)=\prod_{x:X}M_x$ with pointwise
operations (7.1.1). A fraction $m/g^N$ is a pair; $m/g^p=m'/g^{p'}$ is
witnessed by $K$ with $g^K(g^{p'}m-g^pm')=0$. Most results below are
stated with fractions; the localised module $M_g$ is constructed at the
end of this subsection, and 7.1.2 and 7.1.8 are also proved there as
equivalences of modules.

A module $M$ is \emph{weakly quasi-coherent} (7.1.5) if, for every
$g:R$, an element of $M$ given when $g$ is invertible is merely $m/g^N$,
and an $m$ that is zero when $g$ is invertible is merely killed by a
power of $g$. A family is weakly quasi-coherent (7.1.7) if every $M_x$ is.

\paragraph{Boundedness on schemes.} Proposition~3.3.4 extends to a
scheme $X$: if $P:X\to\mathbb N\to\U_l$ is proposition-valued and closed
upwards, and every $P\,x$ merely has an element, then merely one $N$
satisfies $P\,x\,N$ for all $x$. Z-choice gives actual witnesses on each
chart, affine boundedness bounds them there, and finite choice over
the charts gives the bound. No witness is chosen at every point.

\paragraph{Lemma 7.1.2, for fractions.} On a scheme, pointwise merely
existing witnesses of $m/g^p=m'/g^{p'}$ give one exponent for all $x$.
Over a finite open cover, local numerators with a common denominator
whose fractions merely agree on overlaps merely glue to one global $m$,
after raising the denominator by some $g^d$. The argument is first
proved for any family of sets representing fractions, given that raising
numerator and denominator preserves a fraction and that equal fractions
with the same denominator have numerators that agree after finitely many
steps; for modules both are proved.

\paragraph{Lemma 7.1.8.} If $X$ is a scheme and every $M_x$ is weakly
quasi-coherent, then for $g:X\to R$, $M(X)_g\to M(D(g))$ is bijective:
a section over $D(g)$ is merely $m/g^N$ with $m:M(X)$, and an $m$ that
vanishes on $D(g)$ is merely killed by a power of $g$. The first part
is the gluing above; the second applies boundedness to
$P\,x\,K\equiv(g(x)^Km_x=0)$. Applied to $\sum_{x:X}U(x)$, which is a
scheme by Section~5, it gives the form for an open $U\subseteq X$, which
need not be affine.

\paragraph{Sections and images (7.1.10--7.1.12).}
For a constant scalar $g=r$, exchanging the arguments identifies sections
on $D(g)\subseteq X$ with $(\prod_x M_x)^{D(r)}$. Thus the two clauses
of 7.1.8 directly give weak quasi-coherence of $\prod_x M_x$.
We prove 7.1.10 for every scheme $X$: the paper states it for affine $X$,
but applies it in 7.1.12(b) to fibres that need only be schemes.
Inverse image preserves weak quasi-coherence pointwise. Direct image
is the product over the actual fibre, which is a scheme by 5.6.4,
so the strengthened 7.1.10 proves 7.1.12(b).

\paragraph{Finitely presented algebras (7.1.9).}
First, $R$ is weakly quasi-coherent: SQC for the presented localisation
$R_g$ represents functions on $D(g)$, and its fraction and kernel
lemmas give the two clauses. For a finitely presented $R$-algebra $A$,
7.1.10 then makes $(\Spec A\to R)$ weakly quasi-coherent. The canonical
SQC equivalence $A\simeq(\Spec A\to R)$ preserves zero and scalar
multiplication, so the fraction clauses transfer to the underlying
$R$-module of $A$. The bundle statement follows pointwise.
This proof uses the full SAG axioms through 7.1.10; it does not recover
the smaller hypothesis list of the paper's SQC tensor argument.
The general module tensor map in Remark~7.1.3 is not formalised.

\paragraph{Kernels and cokernels (7.1.6, 7.1.13).}
A submodule is a proposition-valued predicate closed under the
operations. The quotient $P/K$ is the effective quotient of $P$ by
$b-a\in K$, with the operations defined on representatives; the
projection is linear and onto, with kernel $K$.

For 7.1.6, let $f:P\to Q$ be linear with $P$, $Q$ weakly quasi-coherent.
A section of $\ker f$ over $D(r)$ is $m/r^N$ with $m\in P$. Then $f(m)$
vanishes on $D(r)$, so $r^Kf(m)=0$, and the section is
$r^Km/r^{K+N}$ with $r^Km\in\ker f$. The zero clause is that of $P$.

Both remaining steps use one lemma: \emph{if $p:P\to Q$ is linear and
onto, and $P$ and $\ker p$ are weakly quasi-coherent, then so is $Q$.}
A section $s$ of $Q$ over $D(r)$ lifts to $P$: the preimages of $s(d)$
form a torsor under $\ker p$, and $D(r)\simeq\Spec R[1/r]$ is affine, so
Theorem~7.3.6 gives a lift $a$. Then $a=m/r^N$ and $s=p(m)/r^N$. If
$p(x)$ vanishes on $D(r)$, then $x\in\ker p$ there, so $x=k/r^N$ with
$k\in\ker p$; $r^Nx-k$ vanishes on $D(r)$, is killed by some $r^K$, and
$r^{K+N}p(x)=0$.

For 7.1.13, apply the lemma to $M\to f(M)$, whose kernel is $\ker f$,
and then to $N\to N/f(M)$, whose kernel is $f(M)$. This avoids the
paper's reduction through embeddings and $M/\ker f$, and its forward
reference to the long exact sequence of 7.3.4.

\paragraph{The localised module.} For a module $M$ over any commutative
ring and an element $r$, $M_r$ is the effective quotient of pairs
$(m,p)$ by $r^K(r^qm-r^pm')=0$ for some $K$. Transitivity uses
$r^q(r^tm-r^pm'')=r^t(r^qm-r^pm')+r^p(r^tm'-r^qm'')$. Addition is
$m/r^p+n/r^t=(r^tm+r^pn)/r^{p+t}$; it respects the relation in the
first argument and is commutative on representatives, which gives the
second argument. The canonical map $M_r\to M^{D(r)}$,
$m/r^p\mapsto(u\mapsto u^pm)$ with $u$ the inverse of $r$, is linear.

\emph{Definition 7.1.5.} This map is an equivalence if and only if the
two clauses used above hold at $r$: surjectivity is the first clause,
and injectivity is the second, since $m/r^p$ maps to $0$ exactly when
$m$ vanishes on $D(r)$. So the clauses are the paper's definition.

\emph{Lemma 7.1.2.} For a scheme $X$, modules $M_x$ and $f:X\to R$,
$M(X)$ is a module over the ring $X\to R$, and
$m/f^p\mapsto(x\mapsto m(x)/f(x)^p)$ is a linear equivalence
$M(X)_f\simeq\prod_{x:X}(M_x)_{f(x)}$. It is injective because pointwise
exponents have a common bound (boundedness on schemes), and surjective
by gluing fractions as above.

\emph{Lemma 7.1.8.} If each $M_x$ is weakly quasi-coherent,
composing with the equivalences $(M_x)_{f(x)}\simeq M_x^{D(f(x))}$ gives
$M(X)_f\simeq M(D(f))$. For an open $U\subseteq X$ and $f:U\to R$, apply
this to the scheme $\sum_{x:X}U(x)$.

\emph{Equalities of modules.} A linear bijection $A\simeq B$ gives $A=B$
in the type of modules. By the record eliminator every module is built
from its fields; by induction on equivalences (univalence) the bijection
may be taken to be the identity; then addition and scaling agree by
linearity, zero and negation are determined by addition, and the laws
are propositions. So 7.1.2 and 7.1.8 hold as equalities of modules over
$X\to R$, as the paper states them.

\emph{Remark 7.1.3.} Only its case $M_r\to M^{D(r)}$ is used, and that is
the map above. The general map $M\otimes A\to M^{\Spec A}$ needs the
tensor product of modules and is not formalised.

\paragraph{Scope and checks.} The proofs take Loc, SQC and Z-choice.
The file \kw{qc-schemes/All.seg} loads Sections~2--5, the torsor
files and the stated parts of Section~7.1 together: 6568 declarations,
no failures, $1.93$ seconds, with 22 expected negative tests rejected.
Tests include positive annihilating exponents, empty fibres and the
nonzero direct image of $R$ along the identity of the point.
The file \kw{cokernel7/All.seg} adds 7.1.6, 7.1.13 and the localised
module, and loads all of Section~7.1: 7167 declarations, no failures,
$3.43$ seconds, with 24 expected negative tests rejected.

\subsection*{Lemma 7.3.5 and Theorem 7.3.6: $H^1$ of an affine scheme vanishes}

\paragraph{Lemma 7.3.5.} Let $M$ be a module over a commutative ring $A$,
and $f_1,\dots,f_n\in A$ with $\sum_k r_kf_k=1$. An element of $M_g$ is
written $x/g^p$. Let $s_{ij}=m_{ij}/(f_if_j)^N$ in $M_{f_if_j}$ with
$s_{jk}-s_{ik}+s_{ij}=0$ in $M_{f_if_jf_k}$. Then there are $p$ and
$V_i\in M$ with
\[
  V_i/f_i^{\,p}-V_j/f_j^{\,p}=s_{ij}\quad\text{in }M_{f_if_j}.
\]
The proof has three steps.
\begin{enumerate}
\item \emph{One exponent.} The cocycle condition says that
  $f_i^Nm_{jk}-f_j^Nm_{ik}+f_k^Nm_{ij}$ is killed by
  $(f_if_jf_k)^{K_{ijk}}$. Let $K$ bound the finitely many $K_{ijk}$ (their
  sum). With $F_i=f_i^{N+K}$ and $\mu_{ij}=(f_if_j)^Km_{ij}$,
  \[F_i\mu_{jk}-F_j\mu_{ik}+F_k\mu_{ij}=0\quad\text{in }M.\]
\item \emph{Powers.} The $F_i$ are again unimodular:
  $\sum_k\rho_kF_k=1$. By recursion on $n$, using that $a+x=1$ gives
  $a(1+x+\dots+x^{q-1})+x^q=1$.
\item \emph{Computation.} With $V_i=\sum_k\rho_k\mu_{ik}$,
  \[F_jV_i-F_iV_j=\textstyle\sum_k\rho_k(F_j\mu_{ik}-F_i\mu_{jk})
    =\sum_k\rho_kF_k\mu_{ij}=\mu_{ij},\]
  which is the statement with $p=N+K$.
\end{enumerate}
The paper writes $s_{ij}=m_{ij}/(f_if_j)$ with one denominator for
each index and uses the identity of step~1 in $M$ itself; steps~1 and~2
are what this needs. The paper's $u_i$ is $-V_i/f_i^{\,p}$, so that
$s_{ij}=u_j-u_i$, while the proof of Theorem~7.3.6 uses
$t_{ij}=u_i-u_j$; the formal statement uses the second convention. Here
elements of $M_g$ are pairs, and equality in $M_g$ is the relation
above; the localised module of Section~7.1 is not needed.

\paragraph{Theorem 7.3.6.} $R$ is a commutative ring, $X$ a type and
$M_x$ an $R$-module for $x:X$. Functions $X\to R$ play the role of $A$
(this is SQC), and $D(g)$ is the set of $x$ where $g(x)$ is invertible.
The hypotheses are:
\begin{itemize}
\item \emph{Loc}: $1\ne0$, and if $a+b$ is invertible, then merely $a$
  or $b$ is;
\item \emph{Z-choice}: if every $B_x$ is merely inhabited, then merely
  there are unimodular $f_1,\dots,f_n:X\to R$ and sections of $B$ over
  each $D(f_k)$;
\item \emph{weak quasi-coherence}, in the form used, Lemma~7.1.8 for
  $U=X$: for $g:X\to R$ the map $M(X)_g\to M(D(g))$ is bijective. On
  elements: a section over $D(g)$ is merely $m/g^N$ for a global $m$,
  and a global $m$ that vanishes on $D(g)$ is merely killed by a power
  of $g$.
\end{itemize}
\emph{Statement:} every $M$-torsor $T$ on $X$ merely has a section, that
is, merely equals the trivial torsor: $H^1(X,M)=0$.

\emph{Proof.}
\begin{itemize}
\item Z-choice gives sections $s_i$ of $T$ over $D(f_i)$, and the
  differences $t_{ij}=s_i-s_j$ are defined on $D(f_if_j)$.
\item By weak quasi-coherence and finite choice, merely
  $t_{ij}=m_{ij}/(f_if_j)^{N_{ij}}$; a common bound $N$ gives one
  exponent.
\item On $D(f_if_jf_k)$ the numerator
  $f_i^Nm_{jk}-f_j^Nm_{ik}+f_k^Nm_{ij}$ is
  $(f_if_jf_k)^N(t_{jk}-t_{ik}+t_{ij})=0$. By weak quasi-coherence and
  finite choice, it is merely killed by powers of $f_if_jf_k$ on all of
  $X$.
\item The hypotheses of Lemma~7.3.5 are now identities at each point
  of $X$. At each $x$ the lemma, for $R$ and $M_x$, gives $p$ and $V_i$;
  with $u_i=V_i/f_i^{\,p}$ on $D(f_i)$, $u_i-u_j=t_{ij}$ on
  $D(f_if_j)$.
\item The corrected sections $(-u_i)\cdot s_i$ agree on overlaps. They
  glue (Lemma~1.2.2): at $x$, the elements equal to all of them form a
  proposition, inhabited because the $D(f_i)$ cover $X$ (unimodularity
  and Loc for finite sums).
\end{itemize}
In the paper's proof, $M(U_i\cap U_j)_{f_if_j}$ should read
$M(X)_{f_if_j}$. The two files have 870 lines.

\paragraph{The paper's hypotheses.} For $X$ affine, Section~7.1 removes
the third hypothesis. The functions $x\mapsto x(f_i)$ on $X\simeq\Spec A$
turn the Z-choice axiom for $A$ into the form above, and Loc gives the
first item. Lemma~7.1.8 gives the third from weak quasi-coherence of
each $M_x$. Hence: if $X$ is affine, each $M_x$ is weakly
quasi-coherent, and Loc, SQC and Z-choice hold, then every $M$-torsor
on $X$ merely has a section. This is the paper's Theorem~7.3.6.

\subsection*{Lemma 7.4.2: cocycles and torsors with local sections}

Let $A$ be an abelian group and $L$ a type with $\tr{L}$. An
$A$-torsor is a set $T$ with a map $A\times T\to T$ such that
$(a+b)\cdot t=a\cdot(b\cdot t)$ and $0\cdot t=t$; for all $t,t'$ there
is exactly one $a$ with $a\cdot t=t'$, written $t'-t$; and $\tr{T}$. An
$L$-cocycle is $c:L\times L\to A$ with $c_{ij}+c_{jk}=c_{ik}$. The
formal statement is an equivalence
\[
  \textstyle\sum_{T:A\text{-torsor}} (L\to T) \;\simeq\; L\text{-cocycles}.
\]
No axiom of~\cite{SAG} is used. Here is the formal proof, written in
the usual style.

\paragraph{From torsors to cocycles.} Given $(T,u)$ with $u:L\to T$, let
$c_{ij}=u_i-u_j$. This is a cocycle, since
$(c_{ij}+c_{jk})\cdot u_k=c_{ij}\cdot(c_{jk}\cdot u_k)=c_{ij}\cdot
u_j=u_i$, and $u_i-u_k$ is the only element with this property.

\paragraph{From cocycles to torsors.} Given $c$, let
\[
  T_c=\{v:L\to A\mid v_m=v_k+c_{km}\text{ for all }k,m\},\qquad
  (a\cdot v)_k=a+v_k,\qquad \hat u_i=(k\mapsto c_{ik}).
\]
We have $\hat u_i\in T_c$ since $c_{im}=c_{ik}+c_{km}$. $T_c$ is a torsor:
it is a set, as a subset of $A^L$; the action laws hold pointwise. For
$v,w\in T_c$, the existence and uniqueness of $w-v$ is a proposition,
so we may assume $k_0:L$. Then $a=w_{k_0}-v_{k_0}$ works:
$a+v_m=(w_{k_0}-v_{k_0})+(v_{k_0}+c_{k_0m})=w_{k_0}+c_{k_0m}=w_m$. It is
unique, since $a'+v_{k_0}=w_{k_0}$ forces $a'=a$. Finally, $\tr{L}$ and
$i\mapsto\hat u_i$ give $\tr{T_c}$.

\paragraph{One composite.} The cocycle of $(T_c,\hat u)$ is $c$: we need
$c_{ij}\cdot\hat u_j=\hat u_i$, that is $c_{ij}+c_{jk}=c_{ik}$ for all $k$.

\paragraph{The other composite.} Let $c$ be the cocycle of $(T,u)$ and
\[
  e:T\to T_c,\qquad e(t)=(k\mapsto t-u_k).
\]
Then $e(t)\in T_c$, since $t-u_m=(t-u_k)+(u_k-u_m)$; $e$ is
equivariant, since $(a\cdot t)-u_k=a+(t-u_k)$; and $e(u_i)=\hat u_i$ by
definition. $e$ is a bijection. It is a proposition that each fibre has
exactly one element, so we may assume $k_0:L$. The element
$y_{k_0}\cdot u_{k_0}$ is mapped to $y$, and $e(t)=y$ gives
$t=(t-u_{k_0})\cdot u_{k_0}=y_{k_0}\cdot u_{k_0}$. By univalence, an
equivariant bijection preserving the points is an equality
$(T,u)=(T_c,\hat u)$: by induction on equivalences it is enough to take
$e$ the identity. Then the two actions are equal by function
extensionality, the torsor laws are propositions, and the points agree.

\paragraph{Comparison with the paper.}
\begin{itemize}
\item The paper takes $T_c=\{u\mid u_i-u_j=c_{ij}\}$ with points
  $k\mapsto c_{ki}$. Its own computation gives the cocycle
  $\tilde c_{ij}=c_{ji}$, so $(T_c,c)$ is a preimage of the transposed
  cocycle $-c$, not of $c$. The definitions above correct this.
\item The paper does not check that $T_c$ is a set with a free and
  transitive action. This is the second place where a point of $L$ is
  assumed.
\item The paper shows that the left side is a set, that the map is an
  embedding (comparing two arbitrary $(T,u)$ and $(T',u')$), and that it
  is surjective. Its map $T\to T'$ depends on a chosen $i:L$. The formal
  proof gives an inverse instead, and compares each $(T,u)$ only with
  $(T_c,\hat u)$. The map $e$ needs no chosen point, and the left side
  being a set is not needed.
\item ``$(T,u)=(T',u')$ is equivalent to an isomorphism'' is one line in
  the paper. In the formal proof it is the only use of univalence, and
  only one direction is needed.
\item The axiom $0\cdot t=t$ is part of the definition, but no step of
  the proof uses it. The left side lies in a larger universe than the
  cocycles.
\end{itemize}
Appendix~\ref{app:torsor} shows how these steps read in the formal text.

\subsection*{Theorem 7.4.7: 2-cocycles and gerbes with local objects}

\paragraph{Definitions.} An $A$-gerbe is a type $G$ that is a 1-type,
merely inhabited and connected ($\tr{x=y}$ for all $x,y$), with for each
$x$ a group isomorphism $\beta_x:(x=x)\to A$, the band. Being a 1-type
follows from the band in~\cite{SAG}; it is included as a field, as in
the note \emph{About torsors and gerbes and liens}, and is a proposition, so the type
of gerbes is the same. The compatibility of the bands with conjugation,
$\beta_y(q^{-1}\ell q)=\beta_x(\ell)$ for $q:x=y$, is not part of the
definition: it holds for any function of $x$, by path induction on $q$.
An $L$-2-cocycle is $c:L^3\to A$ with
$c_{ijk}+c_{ikm}=c_{jkm}+c_{ijm}$, which is the condition of~\cite{SAG}
rearranged. The statement is that
\[
  \textstyle\sum_{G:A\text{-gerbe}}\ \sum_{u:L\to G}\ \prod_{i,j}(u_i=u_j)
  \;\longrightarrow\; L\text{-2-cocycles},\qquad
  c_{ijk}=\beta_{u_i}(p_{ij}\,p_{jk}\,p_{ik}^{-1}),
\]
is an equivalence, for $\tr{L}$. The gerbes are in the same universe as
$L$ and $A$. The paper gives no proof, only a reference to Deligne's
construction~\cite[5.3]{Deligne}.

\paragraph{Torsors.} Paths between two torsors are the same as
equivariant maps: the type of pairs $(Y,f)$ with $f:X\to Y$ equivariant
is contractible, since an equivariant map of torsors is a bijection and
by induction on equivalences. Equivariant maps $X\to X$ are the
translations, so $\mathrm{Aut}(X)\simeq A$. Hence the type of torsors is
a gerbe, the trivial gerbe. It lies one universe higher, but its paths
are small. The image $K$ of the single torsor $A$, given by replacement,
is small, and contains every torsor since all torsors are merely
isomorphic. $K$ is the small trivial gerbe.

\paragraph{Paths in a gerbe.} In a gerbe the paths $x=y$ form an
$A$-torsor, with $a\cdot h=\beta_x^{-1}(a)\,h$. A map of gerbes which
preserves the bands is therefore an equivalence: on paths it is an
equivariant map of torsors, hence a bijection, so the map is an
embedding; and it is surjective because the target is connected.

\paragraph{The cocycle.} $c$ is a 2-cocycle: both sides of the cocycle
law are $\beta_{u_i}(p_{ij}p_{jk}p_{km}p_{im}^{-1})$, after moving
$c_{jkm}$ from $u_j$ to $u_i$ by conjugation with $p_{ij}$.

\paragraph{The proof.} Being an equivalence is a proposition, so we may
assume a point $i_0:L$. We show that each fibre is contractible.
\begin{itemize}
\item \emph{Centre.} Take $K$, with all points equal to the base point
  and as path from $i$ to $j$ the loop of band $c_{i_0ij}$. Its cocycle is
  $c_{i_0ij}+c_{i_0jk}-c_{i_0ik}$, which is $c_{ijk}$ by the cocycle law at
  $(i_0,i,j,k)$.
\item \emph{Contraction.} Given $(G,u,p)$ with cocycle $c$, map $y:G$ to
  the torsor of paths $u_{i_0}=y$, viewed in $K$. This map preserves the
  bands: a loop $\ell$ at $y$ acts on paths by composition, which is the
  translation by $\beta_y(\ell)$. So it is an equivalence. It sends $u_i$
  to the base point, using $h\mapsto\beta(h\,p_{i_0i}^{-1})$, and $p_{ij}$
  to the loop of band $c_{i_0ij}$. A structure identity principle, by
  induction on equivalences, turns this into an equality with the centre.
\end{itemize}

\paragraph{Comparison with Deligne.} Deligne's inverse is the gerbe of
families of torsors $P_i$ with isomorphisms $p_{ij}$ satisfying the
cocycle condition twisted by $c$. The formal proof does not need it:
with a point of $L$, which may be assumed, every gerbe with points is
identified with the trivial one, and the comparison map is Deligne's map
$P\mapsto\mathrm{Hom}(Q_i,P)$ at the single index $i_0$. The proof is about
1400 lines; most of it is the path algebra of the trivialisation.

\subsection*{Theorem 7.4.8: \v{C}ech $H^2$ of a 2-acyclic cover}

\paragraph{Changing the chosen paths.} Let $(G,u,p)$ be a gerbe with
points and paths, with cocycle $c$, and $b$ a 1-cochain. Replacing each
$p_{ij}$ by $b_{ij}\cdot p_{ij}$ (composition with the loop of band
$b_{ij}$) keeps the gerbe and the points, and the new cocycle is
\[
  c_{ijk}+(\delta b)_{ijk},\qquad (\delta b)_{ijk}=b_{ij}+b_{jk}-b_{ik}.
\]
The proof takes place in the torsors of paths: $c_{ijk}$ is the
difference from $p_{ik}$ to $p_{ij}p_{jk}$, and prefixing a path is
equivariant ($q\,(a\cdot h)=a\cdot(q\,h)$, by the compatibility of the
band with conjugation). Hence
$(b_{ij}p_{ij})(b_{jk}p_{jk})=(b_{ij}+b_{jk}+c_{ijk})\cdot p_{ik}$, and
differences in a torsor are unique.

\paragraph{Objects with paths are coboundaries.} For $\tr{L}$,
\[
  \textstyle\sum_{y:G}\prod_i\,(u_i=y)\;\simeq\;
  \sum_{b}\,(c+\delta b=0),
\]
by $(y,r)\mapsto b$, $b_{ij}$ the difference from $p_{ij}$ to
$r_i\,r_j^{-1}$. The paths $r_i$ are part of the data on the left, so
nothing is chosen; a point of $L$ is assumed only to prove that the map
is an equivalence, a proposition. So $c$ is a coboundary exactly when
merely there is an object with paths to all the $u_i$.

\paragraph{Setting.} $X$ is a type with a finite cover $U_1,\dots,U_n$
(every $x$ is merely in some $U_i$) and $F$ a family of abelian groups
over $X$. With $L_x=\sum_i U_i(x)$, a \v{C}ech 2-cocycle is an
$L_x$-2-cocycle for each $x$, as in the paper's proof of Theorem~7.4.5.
For a cocycle $c$, $G_c$ is the family of gerbes given at each $x$ by the
inverse of Theorem~7.4.7. Two cocycles are cohomologous when they differ
by $\delta b$; this is an equivalence relation.

\paragraph{Acyclicity.} The cover is 2-acyclic when $H^1(U_i\cap U_j,F)=0$
for all $i,j$ (including $i=j$) and $H^2(U_i,F)=0$ for all $i$, that is,
every family of torsors over $U_i\cap U_j$ is merely equal to the
trivial family, and every family of gerbes over $U_i$ merely equal to
the trivial family. This is equivalent to: such families merely have a
section. Pointwise, a torsor has an element if and only if it is equal
to $A$, and a gerbe has an object if and only if it is equal to the
trivial gerbe (the trivialisation of Theorem~7.4.7 with one point); for
families both directions are pointwise, by function extensionality.

\paragraph{Finite choice.} Acyclicity gives, for each $i$, merely a
section over the whole $U_i$, and for each $i,j$ merely a section over
the whole $U_i\cap U_j$. Since the cover is finite, these finitely many
mere sections are combined into one by finite choice,
\[
  \textstyle\prod_{i:\mathrm{Fin}\,n}\tr{P_i}\to\tr{\prod_{i:\mathrm{Fin}\,n}P_i},
\]
proved by recursion on $n$ (at higher type). No choice over the points
of $X$ is needed: the sections are given on whole opens by acyclicity.

\paragraph{The theorem.} Let $\check H^2$ be the set of cocycles modulo
being merely cohomologous, and
\[
  \textstyle H^2(X,F)=\tr{\prod_x F_x\text{-gerbe}}_0,
\]
the families of gerbes modulo mere equality; both are effective quotients
constructed from replacement. The map $c\mapsto[G_c]$ induces
$\check H^2\simeq H^2(X,F)$, an equivalence of sets:
\begin{itemize}
\item \emph{well defined}: cohomologous cocycles give the same family,
  by the change of paths and Theorem~7.4.7;
\item \emph{injective}: if $G_c$ and $G_{c'}$ are merely equal, carry
  the points and paths of $c'$ to the gerbes of $c$; acyclicity on the
  $U_i$ gives paths between the two systems of points, conjugating by
  them does not change the cocycle, and the change of paths gives
  $c'=c+\delta b$;
\item \emph{surjective}: for a family of gerbes, acyclicity gives local
  objects on the $U_i$ and paths on the $U_i\cap U_j$, and
  Theorem~7.4.7 gives the cocycle.
\end{itemize}
No separate gluing lemma for 1-types is needed: Theorem~7.4.7 already
constructs the gerbe of a cocycle at each point.

\paragraph{Addition.} The paper states an equality of groups. Addition
on $H^2$ is defined from the sum of gerbes, constructed in three steps.
\begin{itemize}
\item \emph{Sum of torsors.} For $A$-torsors $T,S$, the pairs $(Z,\varphi)$
  with $\varphi:T\times S\to Z$ satisfying
  $\varphi(a\cdot t,s)=a\cdot\varphi(t,s)=\varphi(t,a\cdot s)$ form a
  contractible type (assume points $t_0,s_0$; the centre is $A$ with
  $(t-t_0)+(s-s_0)$). Its centre is $T\otimes S$ with
  $q:T\times S\to T\otimes S$, $q(a\cdot t,b\cdot s)=(a+b)\cdot q(t,s)$,
  and $\mathrm{Hom}(T\otimes S,Z)$ is equivalent to the bilinear maps.
\item \emph{Sum of gerbes.} For gerbes $G_1,G_2$ and $H=G_1\times G_2$, let
  $F_g(h)=(h_1=g_1)\otimes(h_2=g_2)$, a family of torsors on $H$ (loops at
  the source act, so moving the target along a loop $\ell$ is the
  translation by $+\beta(\ell)$). $G_1+G_2$ is the image of
  $g\mapsto F_g$, small by replacement, with $\mu:G_1\times G_2\to G_1+G_2$.
  A loop of $F_g$ is a function $H\to A$; since $H$ is connected and
  inhabited, $c\mapsto(\lambda h.c)$ is an equivalence $A\simeq(H\to A)$,
  and the band is the constant value, with no point chosen. Then
  $\beta(\mu(\ell_1,\ell_2))=\beta(\ell_1)+\beta(\ell_2)$.
\item \emph{Cocycles.} For systems $(u,p)$ and $(u',p')$ with cocycles
  $c,d$, the points $\mu(u_i,u'_i)$ and paths $\mu(p_{ij},p'_{ij})$ have
  cocycle $c+d$. Hence the family of gerbes of $c+d$ is equal to the sum
  of those of $c$ and $d$.
\end{itemize}
The zero cocycle is realised by the trivial gerbe with constant points
and identity paths, so the zero of $H^2$ is the class of the trivial
family. The group laws of $\check H^2$ hold pointwise on cocycles, and
are transferred to $H^2$ through the equivalence, which preserves zero
and addition. So $\check H^2\simeq H^2(X,F)$ as abelian groups.

\paragraph{Refinement.} If a finite cover $V_1,\dots,V_{n'}$ refines
$U_1,\dots,U_n$ through $\rho$ with $V_j\subseteq U_{\rho(j)}$, cocycles
pull back along $(j,v)\mapsto(\rho(j),v)$. Restricting the points and
paths of a gerbe pulls back its cocycle, so the family of gerbes of the
pulled-back cocycle is the same family. Hence, for 2-acyclic covers, the
refinement map on \v{C}ech classes commutes with the two comparison maps
to $H^2(X,F)$; it is therefore an additive isomorphism, independent of
$\rho$, and refinement satisfies the identity and composition laws (by
naturality and injectivity of the comparison maps). Any two 2-acyclic
covers are compared canonically through $H^2(X,F)$. A common refinement
receives the two \v{C}ech refinement maps whether or not it is acyclic
(the pull-back of cocycles and of coboundaries uses no acyclicity); it is
2-acyclicity of the refinement that makes its own comparison map an
isomorphism, and so lets the comparison of the two covers be read through
it. In the formal text the quotient $\check H^2$ is built inside the
segment that assumes acyclicity, so the refinement map on classes is
stated there for 2-acyclic covers.

\section{Summary}

\begin{center}\small
\begin{tabular}{llrrr}
directory & content & hypotheses & lines & time\\\hline
\kw{replacement} & truncation, quotients & --- & 165 & 0.02\,s\\
\kw{sag} & Lemma 2.2.4 & $1\ne0$, Nullst.\ for $R[X]/(f)$ & 807 & 0.03\,s\\
\kw{torsor} & 7.3.5, 7.3.6, 7.4.2, 7.4.7, 7.4.8 & 7.3.6: Loc, Z-choice, WQC & 4922 & 0.97\,s\\
\kw{bounded} & Theorem 3.3.3 & $1\ne0$, two halves of SQC & 428 & 0.07\,s\\
\kw{projective} & Ex.~6.1.9, Lemma 6.2.1 & Loc, SQC instances & 1141 & 0.11\,s\\
\kw{polynomial} & $A[X]$, universal property & --- & 923 & 0.11\,s\\
\kw{presentations} & presentations, Prop.~2.2.2, & Loc, uniform SQC & 1473 & 0.40\,s\\
 & Thm.~3.3.3, Ex.~6.1.9 (algebra) & & & \\
\kw{sag-uniform} & Lemma 2.2.4 & Loc, uniform SQC & 439 & 0.32\,s\\
\kw{tensor} & 3.2.1, 3.2.2, 3.3.1--3.3.3 & 3.2.2, 3.3: uniform SQC; & 825 & 0.33\,s\\
 & & 3.3.3: $1\ne0$ & & \\
\kw{affine} & 3.1, 3.3.4; axioms of 2.1 & Loc, SQC; 3.3.4: Z-choice & 1835 & 0.65\,s\\
\kw{opens} & 2.2.3(b),(c), 4.2 & Loc, SQC; 4.2.7--4.2.9: Z-choice & 686 & 0.89\,s\\
\kw{closed} & 1.3.8, 4.1 & 4.1.2--3: SQC; 4.1.4: all three & 1667 & 0.99\,s\\
\kw{schemes}, \kw{gluing} & 5 (scope above) & Loc, SQC, Z-choice & 2263 & 1.30\,s\\
\kw{bridge7}, & 7.1.1, 7.1.2 (fractions), 7.1.5, & Loc, SQC, Z-choice & 1033 & 1.93\,s\\
\kw{qc-schemes} & 7.1.7--7.1.12; 7.3.6 & & & \\
\kw{cokernel7} & 7.1.2, 7.1.5, 7.1.6, & Loc, SQC, Z-choice & 1266 & 3.43\,s\\
 & 7.1.8, 7.1.13 & & & \\
\kw{horrocks} & \cite{Proj} Lemma 2.1 (rank one) & --- & 5276 & 3.17\,s\\
\kw{laurent} & \cite{Proj} Lemmas 4.1, 4.4 & Loc, SQC & 1217 & 1.17\,s\\
\kw{picard} & \cite{Proj} Cor.~2.3, Prop.~4.6 & Loc, SQC, Z-choice & 2963 & 5.64\,s\\
\kw{pn} & \cite{Proj} Sec.~1, 1.6, 1.7, 7.1--7.3, & Loc, SQC; 7.2, 7.3, & 3017 & 6.61\,s\\
 & Cor.~5.6 & 5.6: Z-choice & & \\
\end{tabular}
\end{center}
Lines count the files of each directory, not the files it shares with
others by links; times are for checking everything a directory imports.
\kw{sag-uniform} combines the proof of \kw{sag} with the
Nullstellensatz of \kw{presentations}.

Loc, SQC and Z-choice are hypotheses, not theorems of the foundational
axioms. For affine $X$, Theorem~7.3.6 needs only these three and weak
quasi-coherence of each $M_x$, as in the paper.
Their interpretation in a model is not formalised here.
Later results beyond those listed above are not formalised.


\appendix
\section{Excerpts of the formal proof of Lemma 7.4.2}\label{app:torsor}

The file is \texttt{primrec/torsor/Torsor.seg}, about 340 lines. The
excerpts below are copied from it. Only \verb|Laws| and \verb|Coc| are
shortened (the file names the intermediate $\Sigma$-types), and
\verb|...| marks an omitted proof. In \sys, \verb|(x : A) * B| is a
$\Sigma$-type, \verb|w.1|, \verb|w.2| are its projections, \verb|Id|
is the identity type, and \verb|calc A { a =< p > b =< q > c }| is an
equational chain. \verb|Tr.Trunc[X].carrier| is the truncation
$\tr{X}$ obtained from replacement, and \verb|univ| is the univalence
parameter.

\paragraph{Torsors, cocycles and differences.} A torsor structure is the
action with its laws. Free and transitive is the contractibility of
$\sum_a a\cdot t=t'$. The difference $t'-t$ (\verb|diff|) is the centre
of contraction, and \verb|uniq| says that it is the only solution.
{\footnotesize
\begin{verbatim}
AD (T : U l) (act : A -> T -> T) : U l =
  (a b : A) (t : T) -> Id T (act (add a b) t) (act a (act b t)) ;
ZE (T : U l) (act : A -> T -> T) : U l = (t : T) -> Id T (act o t) t ;
FT (T : U l) (act : A -> T -> T) : U l =
  (t t' : T) -> isContr ((a : A) * Id T (act a t) t') ;
Laws (T : U l) (act : A -> T -> T) : U l = (sT : isSet T) * Rest T act ;
                                     -- Rest: AD, ZE, FT and ||T||
TorStr (T : U l) : U l = (act : A -> T -> T) * Laws T act ;
Left : U l+ = (T : U l) * (s : TorStr T) * (L -> T) ;
Coc : U l = (c : L -> L -> A) * ((i j k : L) -> Id A (add (c i j) (c j k)) (c i k)) ;

diff (t t' : T) : A = (ft t t').1.1 ;
dact (t t' : T) : Id T (act (diff t t') t) t' = (ft t t').1.2 ;
uniq (t t' : T) (a : A) (p : Id T (act a t) t') : Id A (diff t t') a =
  ap ((b : A) * Id T (act b t) t') A (\ w -> w.1) (ft t t').1 (a, p) ((ft t t').2 (a, p)) ;
\end{verbatim}}

\paragraph{From torsors to cocycles.} The cocycle law is \verb|uniq|
applied to the chain $(c_{ij}+c_{jk})\cdot u_k=c_{ij}\cdot(c_{jk}\cdot
u_k)=c_{ij}\cdot u_j=u_i$ (here \verb|w = (T, s, u)|).
{\footnotesize
\begin{verbatim}
c (i j : L) : A = diff (w.2.2 j) (w.2.2 i) ;
coc (i j k : L) : Id A (add (c i j) (c j k)) (c i k) =
  symH (uniq (w.2.2 k) (w.2.2 i) (add (c i j) (c j k)) (calc w.1 {
    act (add (c i j) (c j k)) (w.2.2 k)
      =< ad (c i j) (c j k) (w.2.2 k) > act (c i j) (act (c j k) (w.2.2 k))
      =< ap w.1 w.1 (act (c i j)) (act (c j k) (w.2.2 k)) (w.2.2 j) (dact (w.2.2 k) (w.2.2 j))
      > act (c i j) (w.2.2 j)
      =< dact (w.2.2 j) (w.2.2 i) > w.2.2 i })) ;
\end{verbatim}}

\paragraph{From cocycles to torsors.} $T_c$, its action and its points:
{\footnotesize
\begin{verbatim}
Cond (v : L -> A) : U l = (k m : L) -> Id A (v m) (add (v k) (c.1 k m)) ;
Tc : U l = (v : L -> A) * Cond v ;
actc (a : A) (v : Tc) : Tc =
  (\ k -> add a (v.1 k),
   \ k m -> transH (ap A A (add a) (v.1 m) (add (v.1 k) (c.1 k m)) (v.2 k m))
                   (symH (assoc a (v.1 k) (c.1 k m)))) ;
uc (i : L) : Tc = (\ k -> c.1 i k, \ k m -> symH (c.2 i k m)) ;
\end{verbatim}}
\noindent Free and transitive is proved by \verb|ftk k0 v w|, given a
point \verb|k0|, and the point is obtained from $\tr{L}$ (\verb|hL|)
because the goal is a proposition:
{\footnotesize
\begin{verbatim}
ft-c : FT Tc actc = \ v w ->
  Tr.Trunc[L].induction @l (\ x -> isContr ((a : A) * Id Tc (actc a v) w))
    (\ x -> isProp-isContr ((a : A) * Id Tc (actc a v) w)) (\ k0 -> ftk k0 v w) hL ;
\end{verbatim}}

\paragraph{The map $e$.} $e(t)=(k\mapsto t-u_k)$, with the proof that it
lies in $T_c$. The proof that $e$ is a bijection (\verb|fib|) uses a
point of $L$ in the same way as \verb|ft-c|.
{\footnotesize
\begin{verbatim}
e (t : T) : PsiM[phi w].Tc =
  (\ k -> diff (u k) t,
   \ k m -> uniq (u m) t (add (diff (u k) t) (c k m)) (calc T {
     act (add (diff (u k) t) (c k m)) (u m)
       =< ad (diff (u k) t) (c k m) (u m) > act (diff (u k) t) (act (c k m) (u m))
       =< ap T T (act (diff (u k) t)) (act (c k m) (u m)) (u k) (dact (u m) (u k))
       > act (diff (u k) t) (u k)
       =< dact (u k) t > t })) ;
eqv : Equiv T PsiM[phi w].Tc =
  (e, qinv T PsiM[phi w].Tc e (\ y -> (fib y).1.1) (\ y -> (fib y).1.2) ...) ;
path : Id Left w (psi (phi w)) =
  sip w.1 w.2.1 w.2.2 PsiM[phi w].Tc eqv PsiM[phi w].str PsiM[phi w].uc (e-act, e-pt) ;
\end{verbatim}}

\paragraph{Univalence.} Induction on equivalences, from univalence, and
its use: an equivariant equivalence preserving the points gives an
identification in \verb|Left|. It is enough to prove this for the
identity (\verb|sip-base|), by function extensionality on the action
and propositionality of the laws.
{\footnotesize
\begin{verbatim}
eqind (T : U l) (P : (T' : U l) -> Equiv T T' -> U l+) (d : P T (idEquiv T))
  (T' : U l) (e : Equiv T T') : P T' e =
  Prelude[l+].subst (Equiv T T') (P T')
    (Axioms[l].idtoeqv T T' ((univ @l T T').sec e)) e ((univ @l T T').sec-eq e)
    (idJ (U l) T (\ X q -> P X (Axioms[l].idtoeqv T X q)) d T' ((univ @l T T').sec e)) ;

SipMot (T : U l) (s : TorStr T) (u : L -> T) (T' : U l) (e : Equiv T T') : U l+ =
  (s' : TorStr T') (u' : L -> T') -> Eqv T T' s s' u u' e
  -> Id Left (T, s, u) (T', s', u') ;
sip (T : U l) (s : TorStr T) (u : L -> T) (T' : U l) (e : Equiv T T')
  : SipMot T s u T' e = eqind T (SipMot T s u) (sip-base T s u) T' e ;
\end{verbatim}}

\paragraph{The lemma.}
{\footnotesize
\begin{verbatim}
lemma-7-4-2 : Prelude[l+].Equiv Left Coc =
  (phi, Prelude[l+].qinv Left Coc phi psi phipsi psiphi) ;
\end{verbatim}}

\begin{thebibliography}{9}
\bibitem{SAG} F.~Cherubini, T.~Coquand, M.~Hutzler.
  \emph{A foundation for synthetic algebraic geometry}.
  arXiv:2307.00073, 2023.
\bibitem{Proj} F.~Cherubini, T.~Coquand, M.~Ritter, D.~W\"arn.
  \emph{Projective space in synthetic algebraic geometry}. 2026.
\bibitem{Deligne} P.~Deligne. \emph{Le symbole mod\'er\'e}.
  Publ.\ Math.\ IH\'ES 73, 147--181, 1991.
\bibitem{CavCoq} E.~Cavallo, T.~Coquand.
  \emph{Type-theoretic replacement and univalent completion: applications
  and interpretations}. TYPES 2025, LIPIcs.
\end{thebibliography}

\end{document}
